\phantomsection
\chapter*{Study Tour}
\addcontentsline{toc}{section}{Study Tour --- Avis Alfarisy}
\penulis{Avis Alfarisy}

\awalcerpen{T}{epatnya hari Jum'at} sekitar jam 11 san, Hari yang dimana sudah saya tunggu-tunggu . Hari itu adalah hari dimana saya dan teman-teman kelas saya akan mengikuti kegiatan study tour ke Jakarta. Saat di umumkan oleh guru ke murid, semua murid sudah sangat antusias karena kegiatan ini bukan hanya untuk belajar, tetapi juga untuk jalan-jalan dan bermain bersama teman-teman.


Malam nya sebelum berangkat, saya sudah menyiapkan semua barang yang akan saya bawa ke sana. Saya mengecek tas berkali-kali supaya tidak ada barang yang ketinggalan di rumah. Saya membawa baju ganti untuk nanti bermain wahana, lalu membawa makanan ringan, minuman, handphone, dan beberapa barang lainnya.


saya merasa senang karena ini adalah kesempatan saya untuk membuat kenangan baru bersama teman-teman saya disana. saya juga tidak sabar untuk pergi ke Dufan karena selama ini saya hanya melihat keseruannya dari sosmed.


Saat pagi sekitar jam 7 an,saya langsung bersiap-siap, saya langsung berangkat ke makam eril. Saat sampai di makam eril, ternyata sudah banyak teman-teman saya yang sudah datang.


Suasana di makam eril sangat rame. ada yang berfoto disana , ada yang membicarakan wahana yang ingin dicoba di Dufan, dan ada juga yang sibuk memeriksa barang bawaan mereka.


Saat sedang melihat teman-teman, saya melihat teman saya yang bernama Daffa datang membawa tas.


Avis : “hey daf”


Daffa : “ey vis, akhirnya berangkat juga ya kita ke jakarta”


Avis : “Iya daf, saya seneng banget”


Daffa : “sama, Saya juga sudah ngga sabar pengen coba wahana yang ada di dufan”


Avis : “Apa saja yang mau kamu coba daf”


Daffa : “nanti vis kita liat saja disana,”


Ga lama setelah ituh , Bu Nunung dan Pak Juju mulai memanggil semua murid untuk berkumpul.


Bu nunung : “semuanya kumpul dulu ya. Kita akan berangkat sebentar lagi,”


Kami semua langsung berkumpul di depan makan eril. Bu Nunung ngecek jumlah murid satu per satu agar tidak ada yang tertinggal di makam eril.


Bu nunung : “pastikan teman di sebelah kalian sudah ada semua. jangan sampai ada teman kalian yang tertinggal disini,”


Murid : “oke buu”


Setelah semua selesai diperiksa, kita mulai naik ke bus kita. saya duduk bersama teman saya yang bernama Daffa di bagian belakang bus. Kita memilih tempat duduk yang dibelakang agar kita bisa bercanda canda di belakang bersama teman teman yang lain.


Saat bus mulai berangkat, semua murid senang termasuk saya . saya melambaikan tangan saya ke orang tua saya


Di dalam bus suasananya sangat seru. Ada yang karoke, ada yang bercanda,dan ada juga yang mulai tidur.


Saya dan teman saya Daffa mengobrol dan bercanda canda di belakang.


Daffa : “Vis, kamu pernah ke Jakarta belum?”


Avis : “belum pernah daf baru kali ini ,”


Daffa : “Oke vis”


Perjalanan menuju Jakarta lumayan lama, tapi kami tidak merasa bosan karena saya bercanda dan karokean di belakang.


Setelah beberapa jam perjalanan, akhirnya kami sampai di Jakarta.


Sebelum bermain ke Dufan, saya dan teman teman saya mengikuti kegiatan belajar dulu . semua murid mengunjungi beberapa tempat dan mendapatkan banyak informasi baru. Bu Nunung menjelaskan berbagai hal kepada murid.


Setelah kegiatan belajar selesai, kita akhirnya menuju tempat yang paling kami tunggu-tunggu, yaitu main wahana di Dufan.


Saat bus sampai di dufan , semua murid langsung senang.


Daffa : “sampe juga akhirnya vis”


Avis : “Alhamdulillah daf ah”


saya melihat banyak wahana yang pengen saya coba


Avis : “Daf, Bingung ey mau naik wahana apa”


Daffa : “sama vis keliatan nya enak ini semua wahana nya”


Kita pun masuk ke area Dufan bersama teman-teman. Kita mulai mencoba berbagai wahana.


Wahana pertama yang kita coba adalah kora-kora.


Avis : “Daf, aman kan ie?”


Daffa : “Aman lah, vis santai weh”


Saat wahana sudah mulai, saya langsung berpegangan erat karna takut.


Avis : “Daf, knp mkin tinggi daf ey?”


Daffa : “santai Vis amann”


Setelah permainan selesai,kita langsung ketwa ketawa.


Salah satu momen yang paling lucu adalah saat Daffa mencoba wahana yang membuat rambutnya berantakan karna rambutnya mirip seperti mie rambut daffa agak kribo.


Ngga terasa waktu nya sudah mulai sore . Bu Nunung kemudian memanggil semua murid untuk berkumpul.


Bu nunung : “ waktunya kembali ke bus. Cek lagi barang kalian sebelum pulang”


Semua murid pun masuk ke bus dan mulai perjalanan pulang.


Di perjalanan pulang, semua murid carape dan semua tertidur


Akhirnya bus sampai kembali di makam eril sekitar jam 8 an kita sampe dengan selamat.


Setelah saya sampe kerumah saya ingin mengeluarkan dompet saya yang ada di saku setelah saya mengecek saku ternyata dompet saya ilang.


Posisi disitu sudah malam dan saya tidak bisa pergi ke makam eril karna sudah terlalu kemalaman dan saya pun memutuskan untuk langsung tidur.


Besoknya tepatnya hari minggu jam 10 saya langsung pergi ke makam eril untuk mengecek dompet saya yaang hilang setelah saya cari-cari ternyata dompetnya tidak ketemu lalu saya memutuskan untuk pulang.


Sesampainya dirumah saya langsung diam dan disitu saya lanjut membantu pekerjaan rumah dan meng ikhlaskan dompet saya, ketika saya lagi menyapu di halaman depan rumah saya, ada orang yang meng hampiri saya dan dia memberikan dompet saya yang ilang dan ternyata orang itu adalah teman saya yang bernama hilal dan saya langsung berterima kasih kepada dia


\textbf{\# TERIMAKASIHTEMAN}

\textbf{\# TEMANSELALUDIHATI}

\jedahias
